\section{Current Law for the Older Roman Books}

\begin{studybox}{Legal scope}
This chapter applies the 1983 Code of Canon Law and universal Latin-Church liturgical law, with publicly located special provisions, through 13 July 2026. It does not apply the Eastern Code to Eastern Catholics, inventory every diocesan decree, or infer a private rescript's terms. No act of Leo XIV repealing or replacing \emph{Traditionis custodes} was located through the cutoff. The controlling general framework therefore remains the 2021 motu proprio, its 2021 Responsa, and the 2023 rescript, subject to valid particular, proper, or special law.
\end{studybox}

The Church's judgment about a Missal's sacramental goodness is not identical with the law governing who may use it. A rite can be venerable and valid while its public celebration is restricted. A priest can possess valid Orders without the faculty, office, or permission needed for a particular act. A religious institute can be fully canonical without enjoying an unrestricted liturgical privilege. These distinctions are indispensable to reading the last half-century.

\subsection{From concession to broad faculty and back to restriction}

\begin{longtable}{>{\bfseries\raggedright\arraybackslash}p{0.15\linewidth}>{\raggedright\arraybackslash}p{0.28\linewidth}>{\raggedright\arraybackslash}p{0.49\linewidth}}
\toprule
\textbf{Date} & \textbf{Act} & \textbf{Legal significance}\\
\midrule
\endhead
1971 & England and Wales concession & Authorized limited celebrations with the older Order under stated conditions. Its famous petitioners did not create a universal right, and the concession did not govern other territories.\\
1974 & Congregation notification & Presented the revised Missal as the book to be used after publication of approved vernacular editions, apart from concessions. It marks practical legal replacement even though later debate concerned formal abrogation.\\
1984 & \emph{Quattuor abhinc annos} & Authorized diocesan bishops to grant an indult under restrictive conditions, including acceptance of the reformed Missal's legitimacy.\\
1988 & John Paul II, \emph{Ecclesia Dei adflicta} & After Archbishop Lefebvre's unlawful consecrations, urged wide and generous application for faithful attached to earlier forms and established the Pontifical Commission Ecclesia Dei. New communities reconciled or formed under this settlement.\\
2007 & Benedict XVI, \emph{Summorum Pontificum} & Described one Roman Rite in ordinary and extraordinary forms, declared the 1962 Missal never juridically abrogated, and gave Latin priests a broad faculty for private Mass and regulated group requests.\\
2011 & \emph{Universae Ecclesiae} & Interpreted the 2007 law, protected access, and assigned supervision to the Ecclesia Dei commission.\\
2019 & Francis, apostolic letter on the Ecclesia Dei commission & Suppressed the former commission and assigned its tasks to a special section of the Congregation for the Doctrine of the Faith; institutional supervision later changed again.\\
16 Jul 2021 & Francis, \emph{Traditionis custodes} & Replaced the 2007 framework, called the books of Paul VI and John Paul II the unique expression of the Roman Rite's lex orandi, restored restrictive episcopal authorization, and abrogated inconsistent prior norms.\\
18 Dec 2021 & Dicastery Responsa & Gave binding answers and explanatory notes on locations, sacraments, readings, bination, and priestly authorization.\\
11 Feb 2022 & Particular decree for FSSP & Granted members a personal faculty to use the 1962 Missal, Ritual, Pontifical, and Breviary under the decree's territorial conditions; it is not a general privilege for every older-book community.\\
20 Feb 2023 & Papal audience rescript & Reserved to the Apostolic See dispensations concerning use of parish churches, erection of personal parishes, and authorization of priests ordained after the motu proprio.\\
2 Jul 2026 & DDF SSPX explanatory note & Following unlawful episcopal consecrations, identified SSPX ministers as schismatic and excommunicated under canon 1364 \S1; declared their sacraments illicit and their Penance and assisted marriages invalid; exhorted faithful to abstain from their events.\\
\bottomrule
\end{longtable}

This sequence resists two mirror myths. The 1962 Missal did not remain universally available on unchanged 1570 terms after 1970. The Church also did not judge it intrinsically harmful or sacramentally false: successive popes conceded, expanded, organized, and then restricted its celebration. Law changed because legislators made prudential and ecclesiological judgments, not because Eucharistic matter and form oscillated.

\subsection{What \emph{Traditionis custodes} presently requires}

Article 2 calls the diocesan bishop the moderator, promoter, and guardian of the liturgical life in his diocese and gives him exclusive competence to authorize use of the 1962 Missal according to the Apostolic See's guidelines. For groups already celebrating, article 3 requires him to determine that they do not deny the validity and legitimacy of the reform, designate locations, establish days, ensure readings are proclaimed in an approved vernacular translation, appoint a suitable delegate, review erected personal parishes, and refrain from authorizing new groups.

The ordinary location may not be a parish church and a new personal parish may not be erected without the reserved dispensation. The 2023 rescript confirms that these are matters reserved to the Apostolic See rather than dispensations a diocesan bishop may silently grant on his own. A concrete diocesan decree remains necessary to know which church, schedule, celebrant, and conditions actually apply.

Article 4 distinguishes priests ordained after 16 July 2021: they must submit a formal request to the bishop, who is to consult the Apostolic See before granting authorization. The 2023 rescript treats this authorization as reserved. Priests already celebrating were also to seek the diocesan bishop's authorization to continue. A priest's general faculty to celebrate Mass is not by itself the special authorization to choose the 1962 book.

Article 6 places institutes of consecrated life and societies of apostolic life erected under the former Ecclesia Dei commission under the competence of the Dicastery for Institutes of Consecrated Life and Societies of Apostolic Life. This establishes governance; it does not publish one identical liturgical faculty for all of them. Article 7 assigns the competent dicasteries supervisory authority. Article 8 abrogates prior norms, instructions, permissions, and customs that do not conform to the motu proprio.

\subsection{Other sacraments and books}

The 2021 Responsa refuses a casual inference from permission for the Missal to permission for every old book. It allows use of the last typical edition of the \emph{Rituale Romanum} only for canonically erected personal parishes under the stated authorization and does not authorize use of the former \emph{Pontificale Romanum} as a general diocesan option. The bishop must also ensure that approved vernacular Scripture is proclaimed, and the permitted older celebration may not be used to construct a hybrid lectionary or private translation.

These general restrictions can yield to valid special law. The FSSP decree dated 11 February 2022 expressly permits its members to celebrate Mass, sacraments, rites, and the Office according to the typical editions in force in 1962, naming the Missal, Ritual, Pontifical, and Breviary. In the fraternity's own churches and oratories the faculty is personal; elsewhere it operates only with the local ordinary's consent, apart from private Mass. The decree also asks the members to consider the provisions of \emph{Traditionis custodes} insofar as possible. Its wording cannot be generalized to the ICKSP, IBP, a monastery, a diocesan priest, or a lay group without their own competent act.

\canonwarning{Canonical erection answers “what is this body in the Church?” Liturgical faculty answers “which books may this minister use, where, and under what conditions?” A pontifical-right institute can still need local consent; a diocesan association can possess a particular authorization; a former Ecclesia Dei relationship does not prove a current unbounded faculty. The community census therefore reports status and public liturgical basis in separate columns.}

\subsection{The bishop's power and its boundary}

A bishop is not a branch manager executing taste. He possesses ordinary, proper, and immediate authority in his diocese under universal law and must guard unity, rights, doctrine, sacraments, and pastoral good. Under current older-book law he assesses groups and celebrants, issues concrete authorizations, ensures competent pastoral care, and applies reserved procedures. He cannot dispense from what the Apostolic See has reserved to itself unless authorized, nor can a priest evade his judgment by invoking a global Internet interpretation.

At the same time, governance should be juridically intelligible. Restrictions affect faithful whose attachment may be doctrinally sound and spiritually fruitful. Clear decrees, stated terms, stable transition, access to sacraments, and charitable dialogue serve communion better than collective suspicion. \emph{Traditionis custodes} itself distinguishes groups that deny the reform's validity from those who do not; evidence about a concrete community matters.

\subsection{What the faithful may and may not infer}

A Catholic may love the 1962 Mass, believe that it embodies ritual goods more successfully, attend an authorized celebration, and work lawfully for its preservation. Preference is not rejection of Vatican II or the postconciliar Missal. The faithful may also prefer the reformed Missal and criticize traditionalist subcultures without treating the old liturgy as poisoned.

No lay petition creates a faculty. A published schedule does not prove current authorization. A valid Eucharist does not automatically fulfill every moral obligation without regard to the community that offers it. After the DDF act of 2 July 2026, the faithful should not apply older SSPX statements as present law; the competent dicastery expressly exhorts abstention from the Society's celebrations and activities. Attendance alone is not mechanically formal schism, but knowing adhesion to a separated ecclesiology is another matter.

Likewise, an indult-era foundation is not necessarily unlawful merely because its public authorization is difficult to locate online. Proper law and administrative acts can be nonpublic. The responsible response is to ask the competent diocese or institute, not to pronounce guilt from silence.

\subsection{Common legal category errors}

\begin{fourstable}{Statement}{Verdict}{Correct category}{Why it fails}
“\emph{Quo primum} gives every priest an immutable right to 1570/1962.” & False. & Continuing papal disciplinary competence. & Pius V himself legislated a current book; his successors revised it and can regulate its use. The 1962 book is not the 1570 edition.\\
“Benedict said never abrogated, so the 2007 rules can never change.” & False. & A later legislator can change disciplinary law. & \emph{Traditionis custodes} expressly abrogates inconsistent norms. Benedict's historical-juridical judgment does not freeze his regulatory scheme.\\
“A valid Mass is automatically licit.” & False. & Validity and liceity. & A valid priest can consecrate while acting without permission, outside communion, or against liturgical law.\\
“A canonical institute automatically has unrestricted 1962 faculties everywhere.” & False. & Juridic status, proper law, and territorial consent. & Erection and liturgical authorization answer different questions.\\
“Every old rite is forbidden except the 1962 Roman Missal.” & False. & Rite-specific and institute-specific law. & Ancient Latin rites and uses, special law, and proper liturgical books cannot be collapsed into one Roman permission.\\
\end{fourstable}

\governingjudgment{The present general law is restrictive, not annihilative. It treats the postconciliar books as the Roman Rite's unique normative expression while allowing bounded celebration of the 1962 Missal through competent authorization and special law. Catholics should neither pretend the 2007 regime still governs nor transform a mutable discipline into a doctrinal condemnation of the old Mass.}
